% ============================================================
% Section 7.2.2  –  n = 30 Instances  (j30, |K| = 4)
% Source: EXP2/exp2_j30_k4.csv  /  exp2_j30.py
% ============================================================
\subsection{$n=30$ Instances}
\label{subsec:n30_results}

\subsubsection*{Experiment Configuration}

Both approaches in this experiment share an identical setup to
ensure a fair head-to-head comparison. The same SSGS warm-start
schedule $S_{\mathrm{SSGS}}$ is passed as the initial solution
to both methods, so any difference in solution quality reflects
the R\&F decomposition strategy alone. The experiment-specific
parameter values are listed in Table~\ref{tab:exp2_params};
note that $|S|=3$ scenarios and a 1-hour time budget
$T_{\max} = 3{,}600\,\text{s}$ are used here, which differ
from the general parameter table (Table~\ref{tab:params_exp})
and are chosen to keep the total experiment runtime tractable
across all 50 instances.

\begin{table}[ht]
\centering
\renewcommand{\arraystretch}{1.35}
\caption{EXP2 parameter settings (j30, $|K|=4$).
         Parameters that differ from the general setup
         (Table~\ref{tab:params_exp}) are marked with~$\dagger$.}
\label{tab:exp2_params}
\begin{tabular}{llll}
\hline
\textbf{Category} & \textbf{Parameter} & \textbf{Value} & \textbf{Note} \\
\hline
\multirow{3}{*}{SAA / Solver}
    & $|S|$                        & 3$^\dagger$     & Scenarios per instance \\
    & $T_{\max}$                   & 3{,}600\,s$^\dagger$ & Total time limit \\
    & $T_{\mathrm{model}}$         & 60\,s           & Per-subproblem limit \\
\hline
\multirow{2}{*}{Baseline}
    & $\omega,\,\delta$            & 1,\;0           & Single-activity window \\
    & \texttt{max\_no\_improve}    & 0               & No early stopping \\
\hline
\multirow{4}{*}{Group-based}
    & \texttt{max\_group\_size}    & 10$^\dagger$    & Produces 3--4 groups for j30 \\
    & $\theta_{\mathrm{merge}}$    & 2.5             & Merge threshold \\
    & $\theta_{\mathrm{corr}}$     & 0.95            & Redundancy pruning \\
    & \texttt{max\_no\_improve}    & 2               & Early stopping post-feasibility \\
    & \texttt{stop\_on\_first\_feasible} & \texttt{False} & Continue after first feasible \\
\hline
\end{tabular}
\end{table}

\noindent
Setting \texttt{max\_group\_size}~$=10$ (rather than 4) allows
the grouper to form larger, fewer groups on j30 instances
($n=30$), typically resulting in 3 groups per instance. This is
appropriate at this scale because each group subproblem remains
tractable within the 60\,s per-subproblem limit while still
providing the solver with a meaningful block of binary variables
to optimise jointly. Setting \texttt{stop\_on\_first\_feasible}~$=
\texttt{False}$ ensures the algorithm continues to all groups
after the first integral solution is found, rather than
terminating immediately, allowing further makespan improvements.

% ── Results ─────────────────────────────────────────────────
\subsubsection*{Aggregate Results}

Table~\ref{tab:j30_k4_summary} summarises the aggregate
performance of both approaches on the 50 PSPLIB j30 instances
with $|K|=4$ resources.

\begin{table}[ht]
\centering
\renewcommand{\arraystretch}{1.35}
\caption{Aggregate results: j30, $|K|=4$, 50 instances}
\label{tab:j30_k4_summary}
\begin{tabular}{lcc}
\hline
\textbf{Metric} & \textbf{Baseline} & \textbf{Group-based} \\
\hline
Feasible instances (of 50)     & 48 (96\%)  & 48 (96\%)  \\
Mean makespan (feasible)        & 99.60      & 98.90      \\
Instances improved over SSGS    & 1 / 48     & 4 / 48     \\
Group-based $<$ Baseline        & \multicolumn{2}{c}{3 / 48 (6\%)}  \\
Both equal                      & \multicolumn{2}{c}{45 / 48 (94\%)} \\
Baseline $<$ Group-based        & \multicolumn{2}{c}{0 / 48 (0\%)}  \\
Mean runtime (feasible)         & 38.4\,s    & 113.9\,s   \\
Max runtime (feasible)          & 60.5\,s    & 259.6\,s   \\
Mean subproblems / groups       & 1.04       & 3.04       \\
\hline
\end{tabular}
\end{table}

\subsubsection*{Feasibility}

Both approaches find a feasible solution for 48 of the 50
instances (96\%). The two exceptions, \texttt{j3045\_10}
and \texttt{j3045\_3}, are the most heavily
resource-constrained instances in the benchmark set and are
not solved by either method within $T_{\max}$.

\subsubsection*{Solution Quality}

On the 48 feasible instances the group-based approach matches
or improves upon the baseline in every case. In 45 of 48
instances (94\%) both methods return the same makespan,
already equal to the SSGS warm-start. In the remaining 3
instances the group-based approach strictly dominates, as
shown in Table~\ref{tab:j30_k4_notable}.

\begin{table}[ht]
\centering
\renewcommand{\arraystretch}{1.35}
\caption{Notable instances: j30, $|K|=4$.
         $z_{\mathrm{SSGS}}$ = SSGS warm-start;
         $z^{\mathrm{base}}$, $z^{\mathrm{grp}}$ = best
         makespan found;
         $\Delta = z^{\mathrm{grp}} - z^{\mathrm{base}}$;
         $t$ = total MIP solve time~(s);
         $m$ = number of groups formed.}
\label{tab:j30_k4_notable}
\begin{tabular}{l r r r r r r r}
\hline
\textbf{Instance} &
$z_{\mathrm{SSGS}}$ &
$z^{\mathrm{base}}$ &
$z^{\mathrm{grp}}$ &
$\Delta$ &
$t^{\mathrm{base}}$ &
$t^{\mathrm{grp}}$ &
$m$ \\
\hline
\multicolumn{8}{l}{\textit{Group-based strictly better}} \\
\texttt{j3037\_1}  & 107 & 107 & \textbf{83}  & $-24$ &  60.2 & 200.5 & 3 \\
\texttt{j3021\_6}  & 109 & 109 & \textbf{100} & $-9$  &  60.2 & 152.7 & 3 \\
\texttt{j306\_5}   & 103 & 103 & \textbf{102} & $-1$  &  60.3 & 259.6 & 5 \\
\hline
\multicolumn{8}{l}{\textit{Both approaches infeasible}} \\
\texttt{j3045\_10} & 139 & \multicolumn{2}{c}{---} & --- & 3{,}623.7 & 2{,}476.6 & 8 \\
\texttt{j3045\_3}  & 148 & \multicolumn{2}{c}{---} & --- & 2{,}556.9 & 1{,}039.4 & 11 \\
\hline
\end{tabular}
\end{table}

The most striking result is \texttt{j3037\_1}, where the
group-based approach reduces the makespan from 107 to 83, a
gain of 24 time units ($-22.4\%$ relative to the SSGS
warm-start) that the baseline entirely misses. A similarly
large improvement is observed on \texttt{j3021\_6} ($-8.3\%$).
These gains are made possible by the combination of
\texttt{stop\_on\_first\_feasible}~$=\texttt{False}$, which
keeps the algorithm running after the first integral solution
is found, and the richer combinatorial context of fixing
an entire group simultaneously, enabling the solver to
reschedule a coherent block of activities jointly
\parencite{pochet2006}. The baseline, constrained to fixing a
single binary variable per subproblem with no early-stopping
awareness of inter-activity interactions, returns the SSGS
warm-start on both instances. Overall, the group-based approach
improves upon the SSGS warm-start on 4 of 48 feasible instances
versus 1 for the baseline, reducing the mean makespan from
99.60 to 98.90.

\subsubsection*{Runtime}

The group-based approach solves on average 3.04 groups per
instance compared to 1.04 subproblems for the baseline. Despite
solving roughly 3$\times$ more subproblems (each containing
multiple binary variables), the mean runtime is 113.9\,s versus
38.4\,s for the baseline, a factor of $2.97\times$. This
overhead reflects the harder subproblems: fixing $|G_g|$ binary
variables simultaneously produces an exponentially larger
branch-and-bound tree than fixing a single variable
\parencite{Wolsey1998}. Nonetheless, all runs complete well
within the 1-hour budget, and the additional solve time is
productive on the three instances where the group-based approach
strictly dominates. The full per-instance results are provided
in Appendix~\ref{app:j30_k4_full}.
